\begin{textsample}{POS Dim 9 – global\_north\_2025\_03 – Score 1.00 – global\_north\_2025\_03/122130311.txt}  \label{ex:f9_pos_302}
Iran rejects Arab plan, calls for one-state solution in Gaza conflict Iran says a two-state solution won't secure Palestinian rights.
IRAN ' S SUPREME Leader Ayatollah Ali Khamenei meets with Yahya Sinwar of Hamas, in Tehran, 2012.
In 2023, Sinwar was asked to defer a large-scale military campaign to give Iran and its allies enough time to prepare for it, but Sinwar chose to go ahead on October 7, says the writer. ( photo credit : Office of the Iranian Supreme Leader/WANA via REUTERS ) Iran has said that it prefers a one-state solution to the conflict in Gaza rather than the recent plan that Arab states have proposed.
The Arab initiative came out of a meeting in Cairo last week.
Iran has now spoken out at a meeting of the Organization of Islamic Cooperation ( OIC ), where foreign ministers from various Muslim countries have gathered.
This is a pivotal time in the region.
There is still a ceasefire in Gaza, and Hamas has rejected Israel ' s attempt to extend @ @ @ @ @ @ @ @ @ @ administration is now open to direct talks with Hamas.
Iran and the Arab states oppose the Trump plan for re-settling Gazans.
The Arab League backs a plan for reconstruction that Egypt has been pushing for Gaza.
Iran ' s Foreign Minister Abbas Araqchi said at the OIC that the two-state solution"will not secure Palestinian rights, emphasizing Tehran ' s support for the establishment of one state representing all the original inhabitants of Palestine."He said,"With due respect to the views of some brotherly countries on the two-state solution, the Islamic Republic of Iran maintains its view that this solution will not lead to the realization of the right of the Palestinian people."This is an important moment for Iran.
It has lost some of its influence in the Middle East after the fall of the Assad regime.
It is trying to claw this back.
It also has less influence in Lebanon, and it wants to use the Houthis to continue to threaten Israel.Leaders of Egypt, Jordan, \textbf{Saudi} Arabia, Qatar @ @ @ @ @ @ @ @ @ @ and the Arab League participate in discussions on the potential displacement of Gazans in Cairo, February 1, 2025. ( credit : screenshot ) No two-state solution"In our view, ' one democratic state ' representing all the original inhabitants of Palestine is the only viable solution,"the Iranian foreign minister said."Given the Israeli regime ' s persistent defiance of the UN Charter, its designation of the UN Secretary-General as persona non grata, the complete obstruction of UNRWA ' s operations, and the unprecedented tragic loss of hundreds of UN staff in Palestine, it is imperative to continue our endeavors for Israeli regime ' s expulsion from the United Nations,"Iran said at the OIC.
Iran ' s state media published the entire speech of the foreign minister, indicating how important they see it.
Iran says that it has"uncompromising support"for the"cause of Palestine."While the Islamic Republic of Iran supports the present resolution, and without prejudice to the foregoing, it would like to put @ @ @ @ @ @ @ @ @ @ outcome of this meeting, details of which will be communicated later with the Secretariat of the OIC,"the Iranians ' said.
Iran says it doesn't want the OIC to be construed as recognizing the"Zionist regime."The Iranian foreign minister went on to say that Iran"earnestly hopes that this auspicious meeting will inspire the international community to take meaningful action to advance justice and peace for the people of Palestine.
May this gathering be a renewed commitment to their honorable cause."Iran is thus trying to claw back some role in the Islamic world.
It is obviously concerned that the \textbf{Saudi} Arabian and Egyptian proposal at the Arab League for Gaza could end with moderation or the replacement of Hamas.
Iran would thus lose out in Gaza as well.

% matched lemmas: saudi
\end{textsample}
